% GENERATED FILE. Edit canonical lessons, then run npm run generate:books.
% canonical-source-hash: fnv1a64:905cbc6c9b1e046b
% canonical-lessons: GU-C336-lekhak, GU-C336-kavi, GU-C336-racvu, GU-C336-atmaktha, GU-C336-kalpana

\chapter{The Author}
\label{ch:gu-the-author}
\hlchaptermodality{336}

\begin{chapteropening}
\textbf{By the end of this chapter:} \emph{I can say who wrote a book and what it is about.}

Complete the last lesson of chapter 336: I can say who wrote a book and what it is about.
\end{chapteropening}

\section[\texorpdfstring{\gu{લેખક} (lekhak)}{લેખક (lekhak)}]{\gu{લેખક} (lekhak) — an author}

\label{lesson:GU-C336-lekhak}



\begin{quote}
Before the new one: say the Gujarati for a proverb, then the Gujarati for to recite.
\end{quote}

\subsection*{You'll want to know: \gu{લેખક}}
\textbf{\gu{લેખક}} — \emph{lekhak} — \textquotedblleft{}an author\textquotedblright{}.

\begin{quote}
\gu{આ} \gu{પુસ્તકનો} \gu{લેખક} \gu{અમારા} \gu{ગામનો} \gu{છે}.
\end{quote}

\begin{quote}
\emph{ā pustakno lekhak amārā gāmno che.}
\end{quote}

\begin{quote}
The author of this book is from our village.
\end{quote}

A masculine or feminine noun, by the person meant.

\begin{cousinweb}
\textbf{\gu{લેખક}} (\emph{lekhak}) is a tatsama, Sanskrit \emph{lekhaka}, \textquotedblleft{}a writer, a scribe\textquotedblright{} (CDIAL 11102), from the same root as \textbf{\gu{લખવું}}, to write.
\end{cousinweb}

\begin{grammarlens}[title={What you've built}]
One more word for telling what happened.
\end{grammarlens}

\subsection*{Your turn}
\begin{itemize}
\raggedright
  \item \emph{Say it:} \emph{lekhak}
  \item \emph{Say it:} \emph{lekhak}, once more
  \item \emph{Say it:} \emph{ā pustakno lekhak amārā gāmno che.} — The author of this book is from our village.
  \item \emph{Recall:} say the Gujarati for a proverb, then the Gujarati for to recite, then say \emph{lekhak} again
\end{itemize}

\begin{tcolorbox}[breakable,colback=teal!4,colframe=teal!35!black,title={Before you move on}]
What does \gu{લેખક} mean? (\textquotedblleft{}An author\textquotedblright{}.) Say it once more.
\end{tcolorbox}
\section[\texorpdfstring{\gu{કવિ} (kavi)}{કવિ (kavi)}]{\gu{કવિ} (kavi) — a poet}

\label{lesson:GU-C336-kavi}



\begin{quote}
Before the new one: say the Gujarati for to recite, then the Gujarati for an author.
\end{quote}

\subsection*{You'll want to know: \gu{કવિ}}
\textbf{\gu{કવિ}} — \emph{kavi} — \textquotedblleft{}a poet\textquotedblright{}.

\begin{quote}
\gu{કવિએ} \gu{સભામાં} \gu{એક} \gu{કવિતા} \gu{વાંચી}.
\end{quote}

\begin{quote}
\emph{kavie sabhāmā̃ ek kavitā vā̃cī.}
\end{quote}

\begin{quote}
The poet read a poem at the gathering.
\end{quote}

A masculine noun.

\begin{cousinweb}
\textbf{\gu{કવિ}} (\emph{kavi}) is a tatsama, Sanskrit \emph{kavi}, \textquotedblleft{}a wise man, a poet\textquotedblright{} (CDIAL 2964). His work is \textbf{\gu{કવિતા}}, a poem.
\end{cousinweb}

\begin{grammarlens}[title={What you've built}]
One more word for telling what happened.
\end{grammarlens}

\subsection*{Your turn}
\begin{itemize}
\raggedright
  \item \emph{Say it:} \emph{kavi}
  \item \emph{Say it:} \emph{kavi}, once more
  \item \emph{Say it:} \emph{kavie sabhāmā̃ ek kavitā vā̃cī.} — The poet read a poem at the gathering.
  \item \emph{Recall:} say the Gujarati for to recite, then the Gujarati for an author, then say \emph{kavi} again
\end{itemize}

\begin{tcolorbox}[breakable,colback=teal!4,colframe=teal!35!black,title={Before you move on}]
What does \gu{કવિ} mean? (\textquotedblleft{}A poet\textquotedblright{}.) Say it once more.
\end{tcolorbox}
\section[\texorpdfstring{\gu{રચવું} (racvũ)}{રચવું (racvũ)}]{\gu{રચવું} (racvũ) — to compose}

\label{lesson:GU-C336-racvu}



\begin{quote}
Before the new one: say the Gujarati for an author, then the Gujarati for a poet.
\end{quote}

\subsection*{You'll want to know: \gu{રચવું}}
\textbf{\gu{રચવું}} — \emph{racvũ} — \textquotedblleft{}to compose\textquotedblright{}.

\begin{quote}
\gu{કવિએ} \gu{નદી} \gu{પર} \gu{એક} \gu{ગીત} \gu{રચ્યું}.
\end{quote}

\begin{quote}
\emph{kavie nadī par ek gīt racyũ.}
\end{quote}

\begin{quote}
The poet composed a song about the river.
\end{quote}

Past, with \gu{તેણે} (he or she …): \gu{તેણે} \gu{ગીત} \gu{રચ્યું}, \emph{teṇe gīt racyũ}.

\begin{cousinweb}
\textbf{\gu{રચવું}} (\emph{racvũ}) is built on Sanskrit \emph{racana}, \textquotedblleft{}arranging, composing\textquotedblright{} (Monier-Williams), from the root \emph{rac}, to fashion: to put words together into a song or a poem.
\end{cousinweb}

\begin{grammarlens}[title={What you've built}]
One more word for telling what happened.
\end{grammarlens}

\subsection*{Your turn}
\begin{itemize}
\raggedright
  \item \emph{Say it:} \emph{racvũ}
  \item \emph{Say it:} \emph{racvũ}, once more
  \item \emph{Say it:} \emph{kavie nadī par ek gīt racyũ.} — The poet composed a song about the river.
  \item \emph{Recall:} say the Gujarati for an author, then the Gujarati for a poet, then say \emph{racvũ} again
\end{itemize}

\begin{tcolorbox}[breakable,colback=teal!4,colframe=teal!35!black,title={Before you move on}]
What does \gu{રચવું} mean? (\textquotedblleft{}To compose\textquotedblright{}.) Say it once more.
\end{tcolorbox}
\section[\texorpdfstring{\gu{આત્મકથા} (ātmakthā)}{આત્મકથા (ātmakthā)}]{\gu{આત્મકથા} (ātmakthā) — an autobiography}

\label{lesson:GU-C336-atmaktha}



\begin{quote}
Before the new one: say the Gujarati for a poet, then the Gujarati for to compose.
\end{quote}

\subsection*{You'll want to know: \gu{આત્મકથા}}
\textbf{\gu{આત્મકથા}} — \emph{ātmakthā} — \textquotedblleft{}an autobiography\textquotedblright{}.

\begin{quote}
\gu{મારા} \gu{દાદાએ} \gu{પોતાની} \gu{આત્મકથા} \gu{લખી}.
\end{quote}

\begin{quote}
\emph{mārā dādāe potānī ātmakthā lakhī.}
\end{quote}

\begin{quote}
My grandfather wrote his autobiography.
\end{quote}

A feminine noun.

\begin{cousinweb}
\textbf{\gu{આત્મકથા}} (\emph{ātmakthā}) is built from Sanskrit \emph{ātman}, \textquotedblleft{}self\textquotedblright{}, and \emph{kathā}, \textquotedblleft{}a story\textquotedblright{} (Monier-Williams): the story of one's own life.
\end{cousinweb}

\begin{grammarlens}[title={What you've built}]
One more word for telling what happened.
\end{grammarlens}

\subsection*{Your turn}
\begin{itemize}
\raggedright
  \item \emph{Say it:} \emph{ātmakthā}
  \item \emph{Say it:} \emph{ātmakthā}, once more
  \item \emph{Say it:} \emph{mārā dādāe potānī ātmakthā lakhī.} — My grandfather wrote his autobiography.
  \item \emph{Recall:} say the Gujarati for a poet, then the Gujarati for to compose, then say \emph{ātmakthā} again
\end{itemize}

\begin{tcolorbox}[breakable,colback=teal!4,colframe=teal!35!black,title={Before you move on}]
What does \gu{આત્મકથા} mean? (\textquotedblleft{}An autobiography\textquotedblright{}.) Say it once more.
\end{tcolorbox}
\section[\texorpdfstring{\gu{કલ્પના} (kalpanā)}{કલ્પના (kalpanā)}]{\gu{કલ્પના} (kalpanā) — imagination}

\label{lesson:GU-C336-kalpana}



\begin{quote}
Before the new one: say the Gujarati for to compose, then the Gujarati for an autobiography.
\end{quote}

\subsection*{You'll want to know: \gu{કલ્પના}}
\textbf{\gu{કલ્પના}} — \emph{kalpanā} — \textquotedblleft{}imagination\textquotedblright{}.

\begin{quote}
\gu{લેખકની} \gu{કલ્પના} \gu{ખૂબ} \gu{સુંદર} \gu{છે}.
\end{quote}

\begin{quote}
\emph{lekhaknī kalpanā khūb sundar che.}
\end{quote}

\begin{quote}
The author has a wonderful imagination.
\end{quote}

A feminine noun.

\begin{cousinweb}
\textbf{\gu{કલ્પના}} (\emph{kalpanā}) is a tatsama, Sanskrit \emph{kalpanā}, \textquotedblleft{}making, forming\textquotedblright{} (Monier-Williams); in Gujarati it is forming pictures in the mind.
\end{cousinweb}

\begin{grammarlens}[title={What you've built}]
One more word for telling what happened.
\end{grammarlens}

\subsection*{Your turn}
\begin{itemize}
\raggedright
  \item \emph{Say it:} \emph{kalpanā}
  \item \emph{Say it:} \emph{kalpanā}, once more
  \item \emph{Say it:} \emph{lekhaknī kalpanā khūb sundar che.} — The author has a wonderful imagination.
  \item \emph{Recall:} say the Gujarati for to compose, then the Gujarati for an autobiography, then say \emph{kalpanā} again
\end{itemize}

\begin{tcolorbox}[breakable,colback=teal!4,colframe=teal!35!black,title={Before you move on}]
What does \gu{કલ્પના} mean? (\textquotedblleft{}Imagination\textquotedblright{}.) Say it once more.
\end{tcolorbox}
